%% ============================================================
%% 片段：s04a-branes   对应原文 PDF 第 83--97 页
%% 本文件由 gw/convert.py 从 arXiv 源码（hep-th/0612073v2）自动生成：
%%   * 公式/编号/引用/图 已经锁定为原文内容，请勿改动任何数学内容与编号；
%%   * 只需把英文正文译为中文，并在每节开头加“本节重点”、关键处加译注。
%% ============================================================


\section{带驯服分歧的几何朗兰兹}

%% ==== 原文 §4：PDF 第 83 页（印刷第 82 页）====
\label{sec:4}

\begin{keytakeaways}
  \begin{itemize}[leftmargin=1.4em,itemsep=2pt,topsep=2pt]
    \item \textbf{物理图像}：把四维 ${\mathcal N}=4$ 超杨–米尔斯理论作 GL 扭曲后置于 $M=\Sigma\times C$，在 $\Sigma$ 远大于 $C$ 的适当极限下（或直接在此拓扑场论中，此时距离无关），理论约化为以 $C$ 上希格斯丛模空间 $\MH$ 为目标的二维西格玛模型；再让 $\Sigma$ 带边界，边界条件由膜标记。$\Psi=\infty$ 给出复结构 $J$ 中的 $B$-模型，$\Psi=0$ 给出辛结构为 $\omega=({\rm Im}\,\tau)\omega_K$ 的 $K$ 型 $A$-模型；$S$-对偶把 $\Psi=\infty$ 映到 $\Psi=0$ 并互换这两类膜，这正是基本的几何朗兰兹对偶。
    \item \textbf{数学工具}：带分歧的希格斯丛模空间 $\MH(\alpha,\beta,\gamma,\eta;p)$ 及其希钦纤维化 $\pi:\MH\to\EUBB$（一般纤维为复环面，$S$-对偶在其上表现为 $T$-对偶）；在复结构 $I$ 下把 $\MH$ 看成抛物丛模空间 $\M(\alpha,p)$ 的余切丛（$\beta=\gamma=0$、$\alpha$ 正则时）或其余切丛的仿射形变，再用典范余迷向膜的 $(\CB_\cc,\CB_\cc)$ 弦环给出微分算子层 ${\mathcal D}_{K_\M^{1/2}}$，并推广到 $K_\M^{1/2}\otimes{\mathcal L}^w$ 型的扭曲。
    \item \textbf{关键结论}：带驯服分歧的几何朗兰兹对偶可表述为 $J$ 型 $B$-膜与 $\M(\alpha,p;G)$ 上相应扭曲 ${\mathcal D}$-模之间的自然对应；扭曲参数 $w(\gamma,\eta)=0\oplus(-\eta-i({\rm Im}\,\tau)\gamma^*)$ 与单值性的本征值 ${}^L\alpha-i{}^L\gamma$ 相等，即“线丛的指数等于单值性的本征值”。这也就是把通常的镜像对称表述改写成 ${\mathcal D}$-模表述的实质。
    \item \textbf{离散数据}：面算符还带有离散电通量 ${\bf e}_0\in\Lambda_{\rm char}/\Lambda_{\rm rt}$ 与离散磁通量 ${\bf m}_0$；存在分歧时，平移 $\alpha$ 可改变 ${\bf m}_0$，对偶地平移 $\eta$ 可改变 ${\bf e}_0$（不改变它们的平移分别落在 $\Lambda_{\rm cort}$ 与 $\Lambda_{\rm rt}$ 中），而 $S$-对偶互换 ${\bf e}_0$ 与 ${\bf m}_0$。
  \end{itemize}
\end{keytakeaways}


\subsection{无分歧情形的回顾}

%% ==== 原文 §4.1：PDF 第 83 页（印刷第 82 页）====
\label{sec:4.1}

从 ${\mathcal N}=4$ 杨–米尔斯理论走向几何朗兰兹的基本步骤见 \cite{ref8}，从该文第 3 节开始。此处把这些步骤简要回顾一遍，因为同一套流程在存在分歧时同样适用。

首先对 ${\mathcal N}=4$ 理论作某种拓扑扭曲（topological twist），即 GL 扭曲。该扭曲依赖于复参数 $t$，给出一族以 $\mathbb{CP}^1$ 参数化的四维拓扑场论（topological field theory）。宜把这个 $\mathbb{CP}^1$ 看成复 $\Psi$ 平面再添上一个无穷远点。$\Psi$ 是 $t$ 与规范耦合参数（gauge coupling）$\tau=\theta/2\pi+4\pi i/e^2$ 的某种组合。${\mathcal N}=4$ 理论中按
$\tau\to (a\tau+b)/(c\tau+d)$ 作用于 $\tau$ 的对偶变换，同样按 $\Psi\to (a\Psi+b)/(c\Psi+d)$ 作用于 $\Psi$。
特别地，对偶操作 $S:\tau\to -1/n_{\mathfrak g}\tau$ 把 $\Psi=\infty$
映到 $\Psi=0$。几何朗兰兹纲领的通常形式正涉及这两个 $\Psi$ 值以及二者之间的对偶。

\begin{yizhu}[原文疑误]
原文此处把对偶操作印作 $S:\tau\to -1/n_{\mathfrak g}\tau$，按字面读是 $-\tau/n_{\mathfrak g}$，即只把 $\tau$ 乘一个常数；但紧接着的同一句话要求该操作把 $\Psi=\infty$ 映到 $\Psi=0$，而常数倍把 $\infty$ 映到 $\infty$，所以唯一自洽的读法是 $S:\tau\to-1/(n_{\mathfrak g}\tau)$（分式线性变换才可能交换 $0$ 与 $\infty$）。第~\ref{sec:4.3}~节对 $\Psi$ 写出的 $S:\Psi\to-1/n_{\mathfrak g}\Psi$ 有同样的排版问题。译文仍照原文保留该式，只在这里提示：凡遇此式须补上分母的括号，否则 $\Psi$ 上的 $\mathrm{SL}(2,\Z)$ 作用无从谈起。
\end{yizhu}

下一步 \cite{ref19,ref20}
是把该理论放到四维流形 $M=\Sigma\times C$ 上，其中 $C$
是希望研究几何朗兰兹纲领的黎曼面（Riemann surface）。
在 $\Sigma$ 远大于 $C$ 的适当极限下，或在拓扑场论中（此时距离无关紧要），紧化后的理论可以用
$\Sigma$ 上的西格玛模型（sigma model）描述，其目标空间是 $C$ 上希格斯丛（Higgs bundle）的模空间（moduli space）$\MH(G,C)$。

在 $\Psi=\infty$ 处，如 \cite{ref8} 第 5 节所述，
所得二维理论是 $\MH$ 在复结构（complex structure）$J$ 中的 $B$-模型。它全纯地依赖于 $\alpha+i\gamma$，并且局部不依赖于 $\beta$ 与 $\eta$。
在 $\Psi=0$ 处，它是辛结构为
% 式 (4.1)
\setcounter{equation}{0}
\begin{equation}\label{eq:4.1}
  \omega=({\rm Im}\,\tau)\omega_K,
\end{equation}

的 $A$-模型，其中因子 ${\rm Im}\,\tau={4\pi/e^2}$ 来自动能项的归一化。
我们称这两种模型分别为 $J$ 型 $B$-模型与 $K$ 型 $A$-模型。
它们各自不依赖于四维规范耦合 $e^2$，因而可以在弱耦合下研究。

\begin{yizhu}[易错点]
关于 $(B,A,A)$ 这类三元组记号与复结构的对应，须注意字母依 $(I,J,K)$ 顺序标记“在哪个复结构上保留 $B$-型超对称”：$B$ 表示保留该复结构的 $B$-模型超对称，$A$ 表示保留相应辛结构的 $A$-模型超对称（原文第~\ref{sec:5}~节给出：$(B,A,A)$ 指保留复结构 $I$ 的 $B$-模型以及复结构 $J$、$K$ 的 $A$-模型的超对称）。据此，本小节所说“$\Psi=0$ 处保留的是 $\omega=(\Im\tau)\omega_K$ 的 $A$-模型”属于 $(B,A,A)$ 型超对称中 $K$ 上的 $A$-型部分；而 $\Psi=\infty$ 处仅靠膜与单值性所能保留的只是复结构 $J$ 上的 $B$-型超对称（即 $(A,B,A)$），要保留完整的 $(B,B,B)$ 还须额外引入威尔逊算符。另一处易错：$(\alpha,\gamma)$ 控制复结构 $J$，故 $B$-模型全纯依赖于 $\alpha+i\gamma$；$\beta$ 控制复结构 $I$；而 $(\beta,\eta)$ 是 $J$ 的 Kähler 参数，故 $B$-模型局部与它们无关——不要把 $\beta,\eta$ 也当成复结构模。
\end{yizhu}

下一步是让 $\Sigma$ 带边界，边界由膜（brane）$\mathcal B$ 标记。$S$-对偶自动把
$J$ 型 $B$-膜与 $K$ 型 $A$-膜互换。这就是基本的几何
朗兰兹对偶。要得到几何朗兰兹对偶的通常表述，还须纳入威尔逊算符（Wilson operator）与 't Hooft 算符
及二者之间的对偶（\cite{ref8} 第 6、8、9、10 节），并把 $K$ 型 $A$-模型
重新解释为 $\M(G,C)$ 上的 ${\mathcal D}$-模（${\mathcal D}$-module）（第 11 节）。

\begin{yizhu}[易错点]
这里其实是两步。第一步是 $S$-对偶：它把规范群 $^LG$ 理论在 $\Psi=\infty$ 的 $J$ 型 $B$-膜与规范群 $G$ 理论在 $\Psi=0$ 的 $K$ 型 $A$-膜对应起来（第~\ref{sec:4.3}~节把两侧的模空间分别写成 $\MH(\dots;{}^LG)$ 与 $\MH(\dots;G)$），所以两侧的目标空间互为对偶的希钦系统，并非同一个 $\MH$；若都当成 $\MH(G)$，就会得出「$\MH$ 自对偶」的错误印象（只在 $G$ 自对偶时才对）。第二步是把 $K$ 型 $A$-模型重新解释为 $\cD$-模（\cite{ref8} 第 11 节）：$A$-膜活在 $\MH$（复维数 $2\dim\M$）上，而 $\cD$-模活在希钦纤维化的底 $\M(G,C)$（复维数 $\dim\M$）上，二者由第~\ref{sec:4.4}~节的 c.c. 膜构造相连。
\end{yizhu}

\subsection{带分歧的西格玛模型}

%% ==== 原文 §4.2：PDF 第 84 页（印刷第 83 页）====
\label{sec:4.2}

下面开始描述在 $C$ 上某点 $p\in C$ 处存在分歧结构时的类比纲领。（为保持记号简单，除确有必要考虑多个分歧点时，我们只讨论单个分歧点的情形。）我们考虑 $M=\Sigma\times C$ 上规范群为 $^LG$ 的规范场论（gauge theory），其中带有支撑在 $\Sigma_p=\Sigma\times p$ 上的“面算符”（surface operator）。对于本文所研究的驯服分歧（tame ramification），取该面算符为熟知类型，由参数
$(\,^L\neg\alpha,^L\neg\beta,^L\neg\gamma,^L\neg\eta)\in\,
^L\TT\times \mathfrak t^\vee\times \mathfrak t^\vee\times \,\TT$ 标记。由于我们把规范群取为 $^LG$，相对于此前的大部分讨论，$^L\TT$ 与
$\TT$ 自然互换。在另行说明之前，取列维子群（Levi subgroup）为 $\mathbb{L}=\TT$。

对偶变换 $S:\tau\to -1/n_{\mathfrak g}\tau$ 把耦合为 $^L\neg \tau$ 的
$^LG$ 理论映为规范群为 $G$、耦合为
$\tau=-1/n_\mathfrak g{}^L\neg\tau$ 的理论。为得到基本的几何朗兰兹对偶，宜取 $^L\tau$ 为纯虚数，并取 $^LG$ 理论的扭曲参数为 $t=i$。
这给出 $\Psi=\infty$ 处的 $B$-模型，并经 $S$ 映为 $t=1$、$\tau$ 为纯虚数的模型。后者给出 $\Psi=0$ 处 $A$-模型的一种便利描述。

\begin{yizhu}[易错点]
注意 $t=i$ 与 $t=1$ 分属对偶两侧：$t=i$ 是映前 $^LG$ 理论（耦合 $^L\tau$）的扭曲参数，经 $S$ 之后才在 $G$ 理论一侧得到 $t=1$。同时 $S$ 把 $\tau$ 换成 $\tau=-1/n_{\mathfrak g}{}^L\tau$，故 $^L\tau$ 纯虚时 $\tau$ 亦纯虚；这两个假设在对偶下互相印证，并非彼此独立。若把 $t=i$ 记到 $G$ 侧，就会以为 $\Psi=0$ 处仍是 $B$-模型，与式~\eqref{eq:4.1} 给出的辛结构 $\omega=({\rm Im}\,\tau)\omega_K$ 及本段所说的 $K$ 型 $A$-模型矛盾。
\end{yizhu}

参数为 $(\,^L\neg\alpha,^L\neg\beta,
^L\neg\gamma,^L\neg\eta)$ 的面算符映为同类型但参数不同的面算符，
其参数为 $(\alpha,\beta,\gamma,\eta)$。这些参数之间的关系已在第~\ref{sec:2.4}~节分析过：

% 式 (4.2)
\setcounter{equation}{1}
\begin{equation}\label{eq:4.2}
  \begin{aligned}(\alpha,\eta)&=(^L\neg\eta,-^L\neg\alpha)\\
                     (\beta,\gamma)&=({\rm Im}\,^L\neg\tau)\,\,(^L\beta^*,^L\neg\gamma^*).\end{aligned}
\end{equation}

第二式即式~\eqref{eq:2.24}，并可如式~\eqref{eq:2.23} 那样反解：

% 式 (4.3)
\setcounter{equation}{2}
\begin{equation}\label{eq:4.3}
  ({}^L\neg\beta,{}^L\neg\gamma)=({\rm Im}\,\tau)(\beta^*,\gamma^*).
\end{equation}

\begin{yizhu}[易错点]
式~\eqref{eq:4.2} 的两行并不对称：第一行是 $(\alpha,\eta)=({}^L\eta,-{}^L\alpha)$，即 $\alpha={}^L\eta$ 无负号、而 $\eta=-{}^L\alpha$ 有负号，这个负号是整行唯一的符号；第二行只是用 $\mathrm{Im}\,{}^L\tau$ 把 $(\beta,\gamma)$ 与 $({}^L\beta^*,{}^L\gamma^*)$ 成比例地联系起来。核对第~\ref{sec:4.3}~节那句 ${}^L\gamma+i{}^L\alpha=({\rm Im}\,\tau)\gamma^*-i\eta$ 时，须同时用式~\eqref{eq:4.3}（得 ${}^L\gamma$）与第一行的负号（得 $-i\eta$）；漏掉负号会把虚部变成 $+i\eta$，从而把复化 Kähler 类的虚部整体反号。
\end{yizhu}

我们的目标是把第~\ref{sec:4.1}~节概述的步骤移植到此情形。
$M=\Sigma\times C$ 上的四维理论在此情形下约化为目标空间为带分歧的希格斯丛模空间的二维西格玛模型。与无分歧时一样，由于目标空间是超 Kähler（hyper-Kähler）流形，该西格玛模型具有 $(4,4)$
超对称。

\begin{yizhu}[易错点]
面算符支撑在 $\Sigma_p=\Sigma\times p$ 上：它是四维理论中的余维 $2$ 算符，沿 $\Sigma$ 铺满二维，沿 $C$ 只是一个点 $p$。因此它在 $C$ 上制造的是 $p$ 处的抛物结构（目标空间随之换成 $\MH(\alpha,\beta,\gamma;p)$、底换成 $\M(\alpha;p)$），并不在 $\Sigma$ 上加任何缺陷，$\Sigma$ 上二维理论的 $(4,4)$ 超对称因此完整保留——依据是目标空间超 Kähler（本段给出的理由），与面算符无关。若把面算符误想成 $\Sigma$ 上的缺陷，就会预期超对称破缺到 $(2,2)$，并把 $\MH(\alpha,\beta,\gamma;p)$ 误当成某种边界条件的模空间。
\end{yizhu}

面参数 $(\alpha,\beta,\gamma)$ 进入模空间的经典几何。其余的面参数 $\eta$ 是第~\ref{sec:2.3}~节比照二维规范场论的 $\theta$ 角引入的。事实上，对秩为 $r$ 的 $G$（且排除 $r$ 与 $g_C$ 都很小的若干例外情形），
$\MH(\alpha, \beta,\gamma;p)$ 的第二 Betti 数为 $r+1$，故该西格玛模型有容纳 $r+1$ 个 $\theta$ 角的余地。其中一个 $\theta$ 角按 \cite{ref8} 式 (4.18) 所述的推理，来自底层四维规范场论的 $\theta$ 角；有效二维理论的另外 $r$ 个 $\theta$ 角则源自 $\eta$。由式~\eqref{eq:4.2} 可见，$S$-对偶把 $\theta$ 角与几何参数互换。

\bigskip\noindent{\it 离散电通量与磁通量}

该西格玛模型还依赖于离散的电通量与磁通量（见 \cite{ref8} 第 7 节，那里采取了略有不同的观点）。我们对此稍作详细描述，否则无法完全精确地说明电磁对偶的作用。

丛 $E\to M=\Sigma\times C$ 具有特征类（characteristic class）
$\xi(E)\in H^2(M,\pi_1(G))$。
记该特征类限制到 $q\times C$ 上为 ${\bf m}_0$，其中 $q$ 为 $\Sigma$ 中的一般点。${\bf m}_0$ 取值于 $H^2(C,\pi_1(G))\cong \pi_1(G)$，
而 $\MH$ 的分支由 ${\bf m}_0$ 标记。

离散电通量需要多作解释。我们对 $G$ 规范场论稍加扭曲，以容许结构群为 $G_{\rm ad}$ 的丛
$E\to M=\Sigma\times C$ 具有如下性质：限制到 $q\times C$ 上（$q$ 为 $\Sigma$ 中的点）时，$E$ 可
提升为 $G$-丛。于是，沿 $\Sigma$ 局部地、沿 $C$ 整体地，理论具有规范群 $G$；但沿 $\Sigma$ 整体上可能存在“扭曲”，
阻碍 $E$ 提升为 $G$-丛。此类丛限制到 $\Sigma\times r$（$r\in C$ 为一般点）上时，具有特征类
${\bf a}_0\in H^2(\Sigma,\pi_1(G_{\rm ad}))$。

在作路径积分（path integral）时，我们引入离散 $\theta$ 角 ${\bf e}_0$，并在路径积分中计入相位因子
% 式 (4.4)
\setcounter{equation}{3}
\begin{equation}\label{eq:4.4}
  \exp(2\pi i\langle {\bf e}_0,{\bf a}_0\rangle).
\end{equation}

于是 ${\bf e}_0$（虽按加法书写）是
$\pi_1(G_{\rm ad})$ 的特征标（character）。然而在 $G$ 规范场论中，自然只考虑在
$\pi_1(G)\subset
\pi_1(G_{\rm ad})$ 上为零的 ${\bf e}_0$（从而使 ${\bf e}_0$ 只对“扭曲”敏感）。于是 ${\bf e}_0$ 是 $\pi_1(G_{\rm
ad})/\pi_1(G)$ 的特征标（这与 $G$ 的中心 ${\mathcal
Z}(G)$ 相同）。等价地，
% 式 (4.5)
\setcounter{equation}{4}
\begin{equation}\label{eq:4.5}
  {\bf e}_0\in \Lambda_{\rm
char}/\Lambda_{\rm rt}.
\end{equation}

若按此方式限制
${\bf e}_0$，则可把 ${\bf a}_0$ 投影到
% 式 (4.6)
\setcounter{equation}{5}
\begin{equation}\label{eq:4.6}
  \pi_1(G_{\rm ad})/\pi_1(G)=\Lambda_\cowt/\Lambda_\cochar.
\end{equation}
对偶变换 $S:\tau\to -1/n_{\mathfrak g}\tau$ 互换
${\bf e}_0$ 与 ${\bf m}_0$。之所以可行，是因为
$\pi_1(\,^L\neg G)=\pi_1(G_\ad)/\pi_1(G)$，且把 $^L\neg G$ 与 $G$
互换后同理。

分歧并不实质改变 ${\bf
m}_0$ 与 ${\bf e}_0$ 的定义；但在有分歧时，它们的作用发生质的变化。正如讨论
式~\eqref{eq:2.6} 时所见，${\bf m}_0$ 可通过平移 $\alpha$ 而改变。对偶地，下面将说明
${\bf e}_0$ 可通过平移 $\eta$ 而改变。$\eta$ 在路径积分中的效应是因子
% 式 (4.7)
\setcounter{equation}{6}
\begin{equation}\label{eq:4.7}
  \exp(2\pi i\langle\eta,\mathfrak{m}\rangle),
\end{equation}
其中 $\mathfrak{m}$ 是把 $E$ 限制到 $\Sigma\times p$ 所得的
$\TT$-丛的特征类（$p\in C$）。在标准 $G$ 规范场论中，$\mathfrak{m}$ 取值于
$\Lambda_{\rm cochar}$；但“扭曲”意味着在目前语境下 $\mathfrak{m}$ 取值于 $\Lambda_{\rm cowt}$。
事实上，
% 式 (4.8)
\setcounter{equation}{7}
\begin{equation}\label{eq:4.8}
  \mathfrak{m} \cong {\bf a}_0~~{\rm
mod}~\Lambda_{\rm cochar}.
\end{equation}
这是关于结构群可约化到环面 $\TT_{\rm ad}$ 的 $G_{\rm ad}$-丛的
二维特征类的一般事实。

由于 $\mathfrak{m}$ 未必取值于 $\Lambda_{\rm cochar}$，通常对 $v\in \Lambda_{\rm char}$ 的对称性 $\eta\to\eta+v$ 未必成立。
确切地说，比较式~\eqref{eq:4.4} 与式~\eqref{eq:4.7} 并利用式~\eqref{eq:4.8} 可知，
$\eta\to \eta+v$ 等价于 ${\bf e}_0\to {\bf e}_0+\bar v$，其中
$\bar v$ 是 $v$ 在 $\Lambda_{\rm char}/\Lambda_{\rm rt}$ 中的像。
不改变 ${\bf e}_0$ 的 $\eta$ 平移恰为按 $\Lambda_{\rm rt}$ 的平移。

这正对应关于 $\alpha$ 的相应陈述：不改变 ${\bf m}_0$ 的平移恰为按 $\Lambda_{\rm cort}$ 的平移。若有多个分歧点
$p_1,\dots,p_s$，这种平行性依然成立。如式~\eqref{eq:3.90} 的讨论所示，不改变 ${\bf m}_0$ 的
$(\alpha_1,\dots,\alpha_s)$ 平移为 $\alpha_i\to\alpha_i+u_i$，其中
$u_i\in \Lambda_{\cowt}$ 且 $\sum_i u_i\in \Lambda_\cort.$
同样，不改变 ${\bf e}_0$ 的 $(\eta_1,\dots,\eta_s)$ 平移为 $\eta_i\to \eta_i+v_i$，其中 $v_i\in \Lambda_\wt$
且 $\sum_iv_i\in \Lambda_\rt$。

\begin{yizhu}[易错点]
离散通量的记账有三处容易混淆。其一，式~\eqref{eq:4.5} 与式~\eqref{eq:4.6} 给出的是同一个有限群的两种写法：由第~\ref{sec:2}~节式~\eqref{eq:2.4}，${\mathcal Z}(G)=\Lambda_{\rm cowt}/\Lambda_{\rm cochar}$，而 $\pi_1(G_{\rm ad})/\pi_1(G)={\mathcal Z}(G)$，故式~\eqref{eq:4.6} 只是把中心的余权格描述与 $\pi_1$ 描述等同起来，式~\eqref{eq:4.5} 则是同一中心的特征标描述（$\Lambda_{\rm char}/\Lambda_{\rm rt}$ 的对偶），不可把两个商当作不同的群。其二，$\mathfrak{m}$ 的取值格：标准 $G$ 规范场论中 $\mathfrak{m}\in\Lambda_{\rm cochar}$，有“扭曲”时 $\mathfrak{m}\in\Lambda_{\rm cowt}$，正因如此 $\eta$ 的周期只是 $\Lambda_{\rm rt}$ 而非整个 $\Lambda_{\rm char}$（对偶地，$\alpha$ 的周期只是 $\Lambda_{\rm cort}$）。其三，多分歧点的陈述“$u_i\in\Lambda_{\rm cowt}$ 而 $\sum_i u_i\in\Lambda_{\rm cort}$”在单点时退化为 $\alpha$ 的单点周期 $\Lambda_{\rm cort}$，初学者常误以为每个点的平移可独立地取 $\Lambda_{\rm cort}$。
\end{yizhu}


\subsection{膜}

 %% ==== 原文 §4.3：PDF 第 87 页（印刷第 86 页）====
\label{sec:4.3}

现在我们把注意力转向两类拓扑场论：$\Psi=\infty$ 处规范群为 $^L\neg G$ 的，以及 $\Psi=0$ 处规范群为 $G$ 的。

在 $\Psi=\infty$ 处得到复结构 $J$ 中的 $B$-模型。在此复结构下，
$\MH(^L\neg\alpha,^L\neg\beta,^L\neg\gamma;p)$ 参数化平坦
$^LG_\C$-丛 $E\to C\backslash p$，其绕点 $p$ 的单值性（monodromy）$V$ 满足一个依赖于 $^L\neg\alpha$ 与
$^L\neg\gamma$ 的条件。

如第~\ref{sec:3.3}~节所给最完整的说明，该条件是：$V$ 在共轭下的轨道，其闭包中含有元素 $U=\exp(-2\pi(^L\neg\alpha-i^L\neg\gamma))$。若 $U$ 正则，这意味着 $V$ 与 $U$ 共轭。若走向另一极端 $U=1$，这意味着 $V$ 是幂幺（unipotent）的。无论如何，对任意 $V\in\, ^L\neg G_\C$，都存在 $^L\neg\alpha$ 与 $^L\neg\gamma$ 的选取（在相差一个外尔变换（Weyl transformation）的意义下唯一），使得以 $V$ 为单值性的平坦联络代表
$\MH(^L\neg\alpha,^L\neg\beta,^L\neg\gamma;p)$ 中的一点。我们只需选取
$^L\neg\alpha$ 与 $^L\neg\gamma$，使 $U$ 落在 $V$ 的轨道闭包中。

\begin{yizhu}[含糊点拨]
「闭包」两字不可省。对 $GL_n$ 及一般约化群，$\overline{\mathcal O_V}\ni1$ 当且仅当 $V$ 的半单部分 $V_s=1$，即 $V$ 幂幺；于是 $U=1$ 的膜恰对应「绕 $p$ 的单值性幂幺」这一分歧几何朗兰兹中最常见的退化情形。反过来，若 $U$ 是正则半单元，它的共轭轨道是闭的，且 $\overline{\mathcal O_V}$ 中元素的半单部分必与 $V_s$ 共轭，故条件退化为 $V$ 半单且与 $U$ 共轭（即正文所说「只意味着 $V$ 与 $U$ 共轭」）。数学上这是因为 $\MH$ 作为抛物希格斯丛模空间只保留轨道闭包层面的信息，故条件是「闭包含 $U$」而非「共轭于 $U$」。
\end{yizhu}

因此，若要借助复结构 $J$ 中的 $B$-模型讨论具有给定单值性的平坦丛，则
$^L\neg\alpha$ 与 $^L\neg\gamma$ 被唯一确定。但
$^L\neg\beta$ 与 $^L\neg\eta$ 是任意的，因为在复结构 $J$ 的观点下它们是 Kähler 参数。复结构 $J$ 中的
$B$-模型当然全纯地依赖于 $^L\neg\gamma+i^L\neg\alpha$，但它局部不依赖于 $^L\neg\beta$ 与 $^L\neg\eta$。整体上，当改变 $^L\neg\beta$ 与 $^L\neg\eta$ 时，$B$-模型带有单值性，我们将在第~\ref{sec:4.5}~节研究。

\begin{yizhu}[易错点]
本节「单值性」（monodromy）一词指两件不同的事，勿混。一处是平坦丛 $E\to C\backslash p$ 绕 $p$ 的单值性 $V$，它决定面算符参数 ${}^L\alpha,{}^L\gamma$ 所满足的条件 $U=\exp(-2\pi({}^L\alpha-i{}^L\gamma))\in\overline{\mathcal O_V}$。另一处是这里的：把 $B$-模型看作参数 $(^L\beta,{}^L\eta)$ 的函数，绕参数空间中的奇点一周时它只回到自身差一个自等价（镜像对称型的 $\mathrm{SL}(2,\Z)$ 单值群），故其物理内容不是参数的单值函数。后者与 $V$ 无关，也不是丛的单值性；第~\ref{sec:4.5}~节讨论的正是它。
\end{yizhu}

在带驯服分歧的几何朗兰兹纲领中，我们从一个具有单值性 $V$ 的平坦 $^L\neg G_\C$-丛
$E\to C\backslash p$ 出发。粗略地说，在
$^L\neg\alpha$ 与 $^L\neg\gamma$ 取到适当值时，$E$ 决定
$\MH(^L\neg\alpha,\,^L\neg\beta,\,^L\neg\gamma;p)$ 上的一个零膜（zerobrane）
${\mathcal B}_E$。（当 $V$ 非正则时，此陈述需要一些补充说明，我们将在第~\ref{sec:4.6}~节给出。）

现在施加 $S$-对偶变换 $S:\Psi\to -1/n_{\mathfrak
g}\Psi$。规范群由 $^LG$ 变为 $G$。$\Psi$
由 $\infty$ 映为 0；所得 $\Psi=0$ 处的模型是具有辛结构 $\omega=({\rm Im\,\tau})\omega_K$ 的 $A$-模型。规范群为 $G$ 的模型的参数
$(\alpha,\beta,\gamma,\eta)$ 可由式~\eqref{eq:4.2} 用 $^L\neg G$ 的参数
$(^L\neg\alpha,\,^L\neg\beta,\,^L\neg\gamma,\,^L\neg\eta)$ 表出。特别地，$\Psi=\infty$ 处的模型所全纯依赖的量
$^L\neg\gamma+i^L\neg\alpha$，按式~\eqref{eq:4.3} 等于 $({\rm Im}\,\tau) \gamma^*-i\eta$。
因此可以预期 $\Psi=0$ 处的 $A$-模型对该变量全纯。事实上，$({\rm Im}\,\tau)\gamma^*-i\eta$ 正是从复结构 $K$ 的观点看的复化 Kähler 类。
（由式~\eqref{eq:3.77}，辛形式 $\omega=({\rm Im}\,\tau)
\omega_K$ 的上同调类正比于 $({\rm Im}\,\tau)\gamma^*$，而 $\eta$ 提供复化 Kähler 类的虚部。）
同样，既然规范群为
$^L\neg G$ 的 $\Psi=\infty$ 处模型局部不依赖于 $^L\neg\beta$ 与 $^L\neg\eta$，就可以预期规范群为 $G$ 的 $\Psi=0$ 处对偶模型局部不依赖于 $\alpha$ 与
$\beta$。这与如下事实相符：作为带有辛结构 $\omega_K$ 的辛簇，$\MH$ 不依赖于 $\alpha$ 与
$\beta$。$A$-模型只依赖这个辛簇及其复化 Kähler 形式 $({\rm Im}\,\tau)\gamma^*-i\eta$。另一方面，在 $A$-模型中无关的复结构 $K$ 则随 $\alpha+i\beta$ 全纯变化。

\begin{yizhu}[含糊点拨]
复化 Kähler 类 $({\rm Im}\,\tau)\gamma^*-i\eta$ 的两项来源可在式~\eqref{eq:3.77} 里逐项读出：$[\omega_J/2\pi]=0\oplus(-\beta^*)$ 与 $[\omega_K/2\pi]=0\oplus(-\gamma^*)$ 的 $\Z$ 分量都是 $0$（带 $e$ 的 $\Z$ 分量只出现在 $[\omega_I/2\pi]=e\oplus(-\alpha^*)$ 中，见式~\eqref{eq:3.76}），故 $({\rm Im}\,\tau)\omega_K$ 的类正比于 $({\rm Im}\,\tau)\gamma^*$；另一项来自 $B$-场，$[B/2\pi]=\eta$。注意 $\gamma$ 与 $\beta$ 的地位关于 $\omega_K$ 与 $\omega_J$ 完全平行（$\gamma$ 配 $\omega_K$、$\beta$ 配 $\omega_J$），这是核对本节全部符号最省力的办法：换言之，$(\beta^*,\gamma^*)$ 就是 $(\omega_J,\omega_K)$ 在 $H^2(\MH)$ 中的坐标，这也正是正文把 $\beta,\gamma$ 与 $\theta$ 角对举、称之为几何参数的原因。
\end{yizhu}

对偶变换 $S$ 把 $J$ 型 $B$-模型中的膜映为 $K$ 型
$A$-模型中的膜。特别地，由带分歧的平坦丛 $E$ 所决定的零膜 ${\mathcal B}_E$ 被映为 $A$-膜 $\hat{\mathcal
B}_E$。我们可以相当具体地看出这将是何种 $A$-膜。无论有无分歧，$\MH$ 都容许希钦纤维化（Hitchin fibration）
$\pi:\MH\to \EUBB$，如第~\ref{sec:3.9}~节所述。希钦纤维化的一般纤维是复环面，在复结构 $I$ 下全纯。按 \cite{ref8} 第 5.4、5.5 节同样的推理，$S$-对偶通过在希钦纤维化的纤维上实现 $T$-对偶而起作用。（这种对偶已从多个数学角度被研究过，包括它在几何朗兰兹纲领中的应用
\cite{ref69,ref70,ref71}。）环面纤维化纤维上的 $T$-对偶把零膜映为绕某条纤维缠绕、并带有平坦 Chan-Paton 线丛 ${\mathcal L}$ 的膜。因此与无分歧时一样，零膜的 $S$-对偶正是这种类型的膜。\cite{ref8} 称这种膜为 $\CMF$ 型膜。

现在可以给出几何朗兰兹对偶的一个初步陈述。陈述时，我们不把 $\MH$ 只看成经典空间，而把它看成定义一个量子西格玛模型，因此在列举其参数时会包含 $\eta$。此外，陈述时把规范群明确写出，并记住电磁对偶互换 $G$ 与
$^L\neg G$。

我们的第一个几何朗兰兹对偶陈述是：对目标为
$\MH(^L\neg\alpha,^L\neg\beta,^L\neg\gamma,^L\neg\eta,p;\,^L\neg
G)$ 的西格玛模型中的每个零膜，都存在目标为
$\MH(\alpha,\beta,\gamma,\eta,p;G)$ 的西格玛模型中相应的 $\CMF$ 型膜；参数照例由
式~\eqref{eq:4.2} 相联系。
更一般地，对目标为
$\MH(^L\neg\alpha,^L\neg\beta,^L\neg\gamma,^L\neg\eta,p;\,^L\neg
G)$ 的西格玛模型中的每个 $J$ 型 $B$-膜（即 $B$-模型在复结构 $J$ 中的每个膜），都存在目标为
$\MH(\alpha,\beta,\gamma,\eta,p; G)$ 的西格玛模型中自然对应的 $K$ 型 $A$-膜。这一对应可扩充为：一侧 $B$-模型的所有结构与另一侧 $A$-模型的所有结构之间的自然对应。

威尔逊算符与 't Hooft 算符的类似物见第~\ref{sec:4.5}~节，一些细节将在第~\ref{sec:4.6}~节澄清。但最迫切的问题是：几何朗兰兹对应通常陈述为左侧是 $J$ 型 $B$-膜（正如上一段），而右侧是适当类型的 ${\mathcal D}$-模，而非 $K$ 型 $A$-膜。为调和这两种观点，我们将采用 \cite{ref8} 第 11 节的做法，读者可参阅该处。

\begin{yizhu}[物理图像补充]
零膜经 $S$-对偶变成“绕希钦纤维缠绕的膜”，其物理机制就是纤维化上的 $T$-对偶：$T^n$ 上的零膜（位置给定、Chan-Paton 丛带平坦联络）在 $T$-对偶下对应于对偶环面上的纤维膜，而平坦联络的平移数据变成对偶侧的威尔逊线。就几何朗兰兹而言，$B$-膜一侧的零膜 ${\mathcal B}_E$ 记录了平坦丛 $E$ 的全部信息（在复结构 $J$ 中它局部化于 $\MH$ 的一个点），而 $A$-膜 $\hat{\mathcal B}_E$ 则在复结构 $I$ 的整条全纯纤维上缠绕，其平坦 Chan-Paton 线丛 ${\mathcal L}$ 的平移数据恰是 $E$ 在希钦纤维上的“位置”的傅里叶对偶。这正是几何朗兰兹中“本征层”图像的起点。把 $\hat{\mathcal B}_E$ 解释为 $\M(\alpha,p)$ 上的 ${\mathcal D}$-模之后，它作为赫克本征层（Hecke eigensheaf）的本征值正是 $E$ 的单值性（即本节末所说的 ${}^L\alpha-i{}^L\gamma$），这就是“指数等于本征值”这一说法的几何来源。
\end{yizhu}


%% ==== 原文 §4.4：PDF 第 89 页（印刷第 88 页）====
\subsection{\texorpdfstring{扭曲 ${\mathcal D}$-模}{扭曲 D-模}}
\label{sec:4.4}


在我们的表述中，几何朗兰兹对应的右侧涉及 $\Psi=0$ 处拓扑场论中的一个膜。
如 \cite{ref8} 一样，到达 $\Psi=0$ 的便利做法是取扭曲参数 $t$ 等于 1，并取
四维 $\theta$ 角 $\theta=2\pi\,{\rm Re}\,\tau$ 为零。
$\Psi=0$ 处的膜是 $K$ 型 $A$-模型的膜。

\begin{yizhu}[易错点]
取 $t=1$ 与取 $\theta=0$ 是两件事，后者在本节约定下等价于 $\mathrm{Re}\,\tau=0$：由 $\tau=\theta/2\pi+4\pi i/e^2$ 得 $\mathrm{Re}\,\tau=\theta/2\pi$，故 $\theta=0$ 当且仅当 $\tau$ 纯虚。这与第~\ref{sec:4.2}~节取 $^L\tau$ 纯虚是同一条件在对偶两侧的表述（$S$ 把纯虚的 $^L\tau$ 变为纯虚的 $\tau$）。正文把 $t=1$ 与 $\theta=0$ 并列为到达 $\Psi=0$ 的做法，可见单靠 $t=1$ 并不足以把 $\Psi$ 钉在 $0$；只有两者同时成立，下面用 $B$-场与复化 Kähler 类所作的论证才严格落在 $A$-模型上。
\end{yizhu}

几何朗兰兹对偶的通常表述涉及
抛物丛（parabolic bundle）模空间 $\M(\alpha,p;G)$ 上的 ${\mathcal D}$-模。
为调和两种表述，我们必须对 $\MH(\alpha,\beta,\gamma,\eta,p;G)$ 上每个 $K$ 型 $A$-膜
配以 $\M(\alpha,p;G)$ 上适当类型的 ${\mathcal
D}$-模。当不至于引起混淆时，为使记号简洁，我们一般省略规范群 $G$ 的标注。

\begin{yizhu}[易错点]
「扭曲」一词在这里指模层取在环层 $\cD_{K_\M^{1/2}\otimes\cL^w}$ 上，其中 $\cL^w$ 象征第一陈类为 $w$ 的线丛，正文把该环层记作 $\cD_{\mathcal L}$，并在本节末把 twist 具体写成 $K_\M^{1/2}\otimes\cL^{-(\eta+i{\rm Im}\,\tau\,\gamma^*)}$。它不同于文献中另一种「twisted $\cD$-模」——把 $\cD$ 沿一个闭 $2$-形式 $\lambda$ 形变所得的 $\lambda$-联络、$\cD_\lambda$-模：这里扭曲的是线丛，算子本身不变，$\lambda$ 这个参数在全文并不出现。两半 twist 的地位也不对等：$K_\M^{1/2}$ 那一半与通常的 $\cD$ 只差 Morita 等价（脚注已指出），真正非平凡的扭曲是 $\cL^w$ 的 $\Lambda_\char$ 分量。
\end{yizhu}

如第~\ref{sec:3.6}~节所述，恰在 $\beta=\gamma=0$ 且
$\alpha$ 正则时，$\MH(\alpha,\beta,\gamma,\eta;p)$ 在复结构 $I$ 下双有理于
$\M(\alpha;p)$ 的余切丛。更确切地说，当 $g_C$ 充分大时，在避开高余维集的意义下，$\MH$ 以余切丛为其稠密开集：

% 式 (4.9)
\setcounter{equation}{8}
\begin{equation}\label{eq:4.9}
  T^*\M(\alpha,p;G)\subset\MH(\alpha,0,0,\eta;p).
\end{equation}
把
$\beta$ 取为零并假定 $\alpha$ 正则，并不是苛刻的限制，因为 $K$ 型 $A$-模型与这些参数无关。但 $A$-模型确实依赖于 $\gamma$，故稍后我们要恢复对 $\gamma$ 的依赖。

此外，取 $\alpha$ 正则表示选取一个包含 $\alpha$ 的仿射外尔房（affine Weyl chamber）。这在某种意义上降低了模型的对称性。
因此可以说，用 ${\mathcal D}$-模来描述对偶（它依赖于这一选取），不如用 $A$-模型与 $B$-模型之间的镜像对称来描述自然。\footnote{也可以基于如下事实作同样的论证：镜像对称保持对偶两侧膜的完整超对称，而用 ${\mathcal D}$-模的描述则不然。关于这一点的解释见第~\ref{sec:5}~节的引言。} 然而，用 ${\mathcal D}$-模的描述是由与数论的类比所推动的，我们希望看清它是如何产生的。

设 $X$ 为任一超 Kähler 流形，复结构为 $I,J,K$，辛形式为
$\omega_I,\omega_J,\omega_K$，且 $[\,\omega_J]=0$
（回忆 $[\,\omega]$ 表示闭形式 $\omega$ 的上同调类）。取正实数 ${\rm Im}\,\tau$，
并考虑 $X$ 上辛结构为 $\omega=({\rm Im}\,\tau)\,\omega_K$ 的 $A$-模型。如 \cite{ref8} 第 11.1 节所述，存在一个特殊的充满空间的 $A$-膜。
其 Chan-Paton 线丛是平凡线丛，带有曲率形式为
$F=({\rm Im}\,\tau)\,\omega_J$ 的联络。
此式保证 $(\omega)^{-1}F$ 是一个复结构，且正是我们称为 $I$ 的那个。这是
\cite{ref72} 意义下余迷向膜（coisotropic brane）的条件。我们称此膜为 c.c. 膜或
典范余迷向膜（canonical coisotropic brane），记作 ${\mathcal B}_{\rm c.c.}$。为使
${\mathcal B}_{\rm c.c.}$ 存在，上同调类
$[\,\omega_J]$ 必须为零，从而平凡线丛才可能具有联络形式
$F=({\rm Im}\,\tau)\,\omega_J$。确实，如第~\ref{sec:3.1}~节末尾最直接说明的，
$[\,\omega_J]=0$ 恰等价于 $\beta=0$。无论如何，为确保 $\MH$ 与
$T^*\M$ 的关系，我们已经取了 $\beta=0$。

\begin{yizhu}[易错点]
这一步可以逐字算出：记 $\omega=({\rm Im}\,\tau)\omega_K$，则 $\omega^{-1}F=({\rm Im}\,\tau)^{-1}\omega_K^{-1}\cdot({\rm Im}\,\tau)\omega_J=\omega_K^{-1}\omega_J$，${\rm Im}\,\tau$ 完全约掉，所得（在超 Kähler 的约定下）正是复结构 $I$。这个抵消不是巧合：它保证 c.c. 膜的构造与 $\tau$ 的具体取值无关，$\tau$ 只以整体标度进入 $\omega$。反过来，若把联络取成 $F=c\,\omega_J$（$c\ne\pm1$），则 $\omega^{-1}F=c\,(\omega_K^{-1}\omega_J)$，平方为 $-c^2\ne-1$，就不再是复结构——可见此处的系数不是约定而是条件。
\end{yizhu}

对任一膜 $\mathcal B$，$(\CB,\CB)$ 弦构成一个环。在西格玛模型微扰论的所有阶，人们可以通过考虑在开集 ${\mathcal U}\subset X$ 内正则的波函数来“局域化”
$(\CB,\CB)$ 弦。在 ${\mathcal U}$ 内正则的开弦构成一个环；令 ${\mathcal U}$ 变化，就得到微扰论所有阶意义下 $X$ 上的一个环层。
一般说来，这种环层的构造只在微扰论所有阶意义上成立，而非严格陈述。
然而，如 \cite{ref8} 所述（数学文献中当然也已知道；更深入的分析见 \cite{ref73}），
在超 Kähler 流形 $X$ 的典范余迷向膜
$\CB_\cc$ 这一情形，可以在一定条件下以一种非常特殊的方式超越微扰论。所需条件是 $X$
应以某个余切丛 $T^*{\mathcal M}$（${\mathcal M}$ 为某个 Kähler 流形）作为稠密开集，
且该等同在复结构 $I$ 下全纯，而 $\omega_K$ 是余切丛自然全纯二形式的虚部。在此情形下，可以对 $X$ 中每个形如 $T^*{\mathcal U}$（${\mathcal U}$ 为
$\M$ 中开集）的开集配以 $(\CB_\cc,\CB_\cc)$ 弦的环。这种对应给出 $\M$ 上的一个环层。该环层恰是
$\M$ 上作用在某个线丛\footnote{更一般地，如 \cite{ref8} 第 11.1 节所回顾，${\mathcal L}$ 可以是线丛的复幂，或这类复幂的张量积。环层 ${\mathcal
D}_{\mathcal L}$ 在把 ${\mathcal L}$ 扭以平坦线
丛时不变，因此在上述更一般的情形下仍有意义
\cite{ref74}。
出于同样的理由，下文不关心
$K_\M^{1/2}$ 作为线丛是否存在或是否唯一。它确实
存在，但不总是唯一。} $\mathcal L$ 的截面上的微分算子层。我们记该环层为 ${\mathcal
D}_{\mathcal L}$。

若 $X$ 的西格玛模型具有时间反演对称性（在其 $\theta$ 角为零时如此），则 ${\mathcal L}$ 可以用特别简单的方式确定。记 $K_\M$ 为 $\M$ 的典范线丛，则时间反演对称性蕴含
${\mathcal L}=K_\M^{1/2}$，而环层为 ${\mathcal D}_{K_\M^{1/2}}$。

要把这一点用于我们的问题，我们希望目标为 $\MH(\alpha,\beta,\gamma,\eta;p)$ 的西格玛模型的 $\theta$ 角全为零。我们已经把四维参数 $\theta$ 取为零；
西格玛模型其余的 $\theta$ 角正是 $\eta$。
故我们通过取 $\eta=0$ 来获得时间反演对称性。我们已经通过取 $\gamma=0$ 简化了问题，因此合并后的条件等价于复化 Kähler 类
$({\rm Im}\,\tau)\gamma^*-i\eta$ 为零。当然，条件 $\eta=0$ 等价于 $^L\neg\alpha=0$。

$(\CB_\cc,\CB_\cc)$ 弦的层与环层
${\mathcal D}_{K_\M^{1/2}}$ 相同，这一事实使得一个非常一般的构造成为可能。
设 $\CB'$ 为 $X$ 上任一 $K$ 型 $A$-膜。
对 $\CB'$ 可配以 $(\CB_\cc,\CB')$ 弦的层。
它自动是 $(\CB_\cc,\CB_\cc)$ 弦的模层，即
${\mathcal D}_{K_\M^{1/2}}$ 的模层。于是我们得到把 ${\mathcal D}_{K_\M^{1/2}}$-模自然配给每个 $K$ 型 $A$-膜的方法。

\begin{yizhu}[含糊点拨]
这个构造与 $S$-对偶无关，对任意超 Kähler $X$ 与任意 $K$ 型 $A$-膜 $\CB'$ 都成立：$(\CB_\cc,\CB')$ 弦层自动是 $(\CB_\cc,\CB_\cc)$ 弦环的模，而后者就是 $\cD_{K_\M^{1/2}}$。所以几何朗兰兹里的「$B$-膜 $\leftrightarrow$ $\cD$-模」是两步的复合：先用第~\ref{sec:4.3}~节的 $S$-对偶把 $\MH(\dots;{}^LG)$ 上的 $J$ 型 $B$-膜送到 $\MH(\dots;G)$ 上的 $K$ 型 $A$-膜，再套用这里的 c.c. 膜构造。前者依赖对偶与 $\Psi=\infty\to\Psi=0$ 这一步，后者只是超 Kähler 几何的定理；这也解释了为何本小节余下的论证全部在 $A$-膜一侧进行。
\end{yizhu}

把它与从 $S$-对偶学到的内容结合起来可见：
$\MH(0,{}^L\neg\beta,0,{}^L\neg\eta,
p;{}^L\neg G)$ 上的 $J$ 型 $B$-膜自然地对应于
$\M(\alpha,p;G)$ 上环层 ${\mathcal D}_{K_\M^{1/2}}$ 的一个模层（照例取 $\alpha=\,^L\neg\eta$）。这本质上\footnote{通常的表述涉及环层
${\mathcal D}$，即作用于函数的微分算子层。该环层 Morita 等价于 ${\mathcal
D}_{K_\M^{1/2}}$。此处不考虑一些微妙之处，它们涉及这一 Morita 等价对 $C$ 上自旋结构选取的依赖。}就是带驯服分歧的几何朗兰兹纲领的通常表述，对应
$^L\neg\alpha=\,^L\neg\gamma=0$ 的情形，换言之，即具有幂幺单值性的平坦
$^L\neg G$ 丛。

与 \cite{ref8} 一样，$(\CB_\cc,\CB_\cc)$ 弦就是 $K_\M^{1/2}$
上的微分算子（而非某个更一般线丛上的微分算子）这一事实，可以更显式地看出，而不必依赖
时间反演对称性。为此我们利用 $\MH$ 上另一个重要的
膜：支撑在 $\varphi=0$ 处、即余切丛 $T^*\M$ 的零截面上，并带有平凡 Chan-Paton 线丛的膜 ${\mathcal B}'$。对
$(\CB_\cc,\CB')$ 弦作量子化表明，它们是
$K_\M^{1/2}$ 的截面。由于它们也给出
$(\CB_\cc,\CB_\cc)$ 弦的模层，后者正是作用在 $K_\M^{1/2}$ 截面上的
微分算子。

\bigskip\noindent{\it 恢复参数依赖}

$\gamma=\eta=0$ 的情形就是 $U=\exp(-2\pi(^L\neg
\alpha-i^L\neg \gamma))$ 等于 1 的情形。若局限于此情形，就意味着只能对幂幺单值性的情形描述几何朗兰兹。要超出这一情形，必须恢复对 $\gamma$ 与 $\eta$ 的依赖。

下文将看到，纳入 $\gamma$ 或 $\eta$ 的效果，是把作用在
$K_\M^{1/2}$ 上的微分算子换成作用在 ${\mathcal S}$ 上的微分算子，其中粗略地说 ${\mathcal S}$ 是一个更一般的线丛——
更确切地说，它是线丛的若干复幂的张量积。作用在这种
${\mathcal S}$ 上的微分算子层有意义，尽管 ${\mathcal S}$ 本身不能定义为线丛。
我们的目标是
{\it (i)} 论证如下断言：对任意 $\gamma$ 与 $\eta$，$\MH$ 上 $K$ 型 $A$-膜都自然地联系于这类环层的模；以及 {\it (ii)} 确定 ${\mathcal S}$。

先论证断言 {\it (i)}。考虑 $\gamma\not=0$、而 $\eta$ 仍为零的效应。
如第~\ref{sec:3.6}~节所述，其结果是 $\MH$ 不再双有理于 $T^*\M$；相反，它以 $T^*\M$ 的一个仿射形变作为稠密开集（这给出与 $\theta\not=0$ 情形类似的情况，如 \cite{ref8} 第 11.3 节所分析）。仅以仿射形变代替余切丛这件事本身，并不改变 c.c. 膜
$\CB_\cc$ 的定义（只要保持 $\beta=0$；见下），也不改变“$(\CB_\cc,\CB_\cc)$ 弦的层是作用在某个‘线丛’ ${\mathcal S}$ 上的微分算子层”这一论证。然而，它破坏了利用时间反演对称性来证明 ${\mathcal S}=K_\M^{1/2}$ 的做法。
事实上，相关时间反演操作的定义要求 $\MH$ 具有对称性 $\varphi\to-\varphi$
（在复结构 $I$ 下全纯地作用）；当 $T^*\M$ 形变为仿射丛时，这一对称性不复存在。另一个用以确定 ${\mathcal S}\cong K_\M^{1/2}$ 的论证用的是支撑在 $\M\subset T^*\M$（即余切丛零截面）上的膜 ${\mathcal B}'$。
当 $\gamma\not=0$ 时这一论证同样失效，因为 $T^*\M$ 的仿射形变不容许全纯截面。

\begin{yizhu}[含糊点拨]
$\varphi\to-\varphi$ 是余切纤维的逐纤维缩放：若 $q^\alpha$ 为 $\M$ 上的坐标、$p_\alpha$ 为纤维坐标，则 $\Omega_X=\sum_\alpha dp_\alpha\wedge dq^\alpha$ 在此变换下变号，$\omega_K={\rm Im}\,\Omega_X$ 也随之变号。它能充当时间反演所需的对称性，是因为它整体有定义并且固定零截面 $\M=\{\varphi=0\}$。一旦 $T^*\M$ 形变为仿射丛，$\varphi$ 只是仿射坐标、$\pi$ 不容许全纯截面，就无所谓「零点」，$\varphi\to-\varphi$ 无从定义——这正是本段说该对称性不复存在的原因，也与紧接的 ${\mathcal B}'$ 论证失效同源（两者都依赖零截面）。
\end{yizhu}

现在考虑 $\eta\not=0$（而 $\gamma$ 为零或非零）的效应。由于对 $({\rm Im}\,\tau)\gamma^*-i\eta$ 的全纯性，$\eta$ 的效应必然与 $\gamma$ 类似，但看清其由来是有意思的。$\eta\not=0$ 导致西格玛模型的 $B$-场（我们简称 $B$）非零。因此，定义典范余迷向膜时需加小心。记该膜的 Chan-Paton 丛 ${\mathcal L}$ 的曲率为 $F$。在存在 $B$-场时，余迷向膜的条件是
$\omega^{-1}(F+B)$ 应为可积复结构，我们将取其为
$I$。为此取 $F+B=({\rm Im}\,\tau)\omega_J$。
我们仍假定 ${\mathcal L}$ 拓扑平凡。这在上同调类层面要求 $[F]=0$，故 $[B]=({\rm Im}\,\tau)[\,\omega_J]$。
由式~\eqref{eq:3.77} 可知 $[\,\omega_J/2\pi]=-\beta^*$，而
$[B/2\pi]=\eta$，故 c.c. 膜只能对
$\beta^*=-({\rm Im}\,\tau)^{-1}\eta$ 定义。因此对非零 $\eta$ 必有 $\beta\not=0$，于是 $\MH$ 再次双有理于 $T^*\M$ 的仿射形变，而非 $T^*\M$ 本身。

此处我们假定 ${\mathcal L}$ 拓扑平凡，
但也可以构造非平凡 ${\mathcal L}$ 的膜。如下文所述，由这种方式得到的结果 Morita 等价于从
${\mathcal L}$ 拓扑平凡情形所能学到的内容。

如前，一旦定义了膜 $\CB_\cc$，标准论证表明
$(\CB_\cc,\CB_\cc)$ 弦的层是作用在某个“线丛”
${\mathcal S}\to \M$ 上的微分算子层。然而，如形变 $\gamma$ 的情形一样，我们无法论证 ${\mathcal S}=K_\M^{1/2}$。也不能用时间反演不变性来证明 ${\mathcal S}=K_\M^{1/2}$，因为 $\eta\not= 0$ 破坏时间反演
对称性。至于尝试用膜 ${\mathcal B}'$ 来论证会得到什么，留待后文描述。

为确定 ${\mathcal S}$，我们如下进行。已知在
$\gamma=\eta=0$ 时 ${\mathcal S}=K_\M^{1/2}$。故写 ${\mathcal S}=
K_\M^{1/2}\otimes {\mathcal S}'$，其中 ${\mathcal S}'$ 在
$\gamma=\eta=0$ 时平凡。作为 $\M$ 上线丛复幂的张量积，${\mathcal S}'$ 有第一陈类 $w\in H^2(\M,\C)$。
（具体地，若 ${\mathcal S}'=\otimes_i{\mathcal L}_i^{g_i}$，其中 ${\mathcal L}_i$ 为普通
线丛、$g_i\in \C$，则
$w=\sum_ig_ic_1({\mathcal L}_i)$。）我们形式上写 ${\mathcal S}'={\mathcal
L}^w$，以表达 ${\mathcal S}'$ 的第一陈类为 $w$ 这一陈述。我们希望确定 $w$。

事实上，$w$ 关于 Kähler 参数 $z=({\rm Im}\,\tau)\gamma^*-i\eta$ 是线性的。这完全可由
\cite{ref8} 第 11.3 节用来证明类似指数 $f(\theta)$ 线性依赖于四维参数 $\theta$ 的同一推理得到。其想法是在弱耦合下计算：取 $e^2=4\pi/{\rm Im}\,\tau$ 很小而 $z$ 固定。
于是 $w$ 可按 $z/{\rm Im}\,\tau$ 作微扰展开计算。具体地，微扰论被用来计算显式描述
$(\CB_\cc,\CB_\cc)$ 弦层结构时出现的余循环。微扰论的结构使得只可能出现 $z$ 的线性项。更多细节见 \cite{ref8} 中式 (11.39) 的讨论。

\begin{yizhu}[含糊点拨]
原文断言 $w$ 关于 $z=({\rm Im}\,\tau)\gamma^*-i\eta$ 线性，并把论证推给 \cite{ref8} 式 (11.39)；被省略的是“高阶项为何不出现”的具体机制。其要点是：$w$ 由 $(\CB_\cc,\CB_\cc)$ 弦环结构的余循环（结合律的障碍类）读出，而 $z$ 只以整体平移闭弦背景（$B$-场与复化 Kähler 类）的方式进入；在 $e^2\to0$、$z$ 固定的弱耦合展开中，该余循环对 $z$ 的依赖受微扰阶数控制，高阶贡献相互抵消，故只留线性项。原文用 \cite{ref8} 第 11.3 节中 $f(\theta)$ 对 $\theta$ 的线性依赖作类比，两者的论证结构相同。核对时应注意：这一步依赖 $\alpha$ 正则与弱耦合展开，原文并未给出强耦合下的证明——这也是作者把整套论证称为启发式的原因。
\end{yizhu}

尚需确定 $w$ 究竟等于 $({\rm Im}\,\tau)\gamma^*-i\eta$ 的哪一个线性函数。
$w$ 取值于 $\M(\alpha,p)$ 的带复系数的第二上同调。
由式~\eqref{eq:3.47} 与式~\eqref{eq:3.56}，这正是 $(\Z\oplus \Lambda_\char)\otimes\C$。
由于 $\eta+i({\rm Im}\,\tau)\gamma^*$
取值于 $\Lambda_\char\otimes \C$，存在自然定义的线性函数

% 式 (4.10)
\setcounter{equation}{9}
\begin{equation}\label{eq:4.10}
  w(\gamma,\eta)=0\oplus (-\eta-i({\rm Im}\,\tau)\gamma^*),
\end{equation}

我们断言这就是正确结果。

\begin{yizhu}[易错点]
${\mathcal D}$-模的 twist 约定有四层含义需要分清。其一，${\mathcal S}=K_\M^{1/2}\otimes{\mathcal L}^w$ 中的 ${\mathcal S}$（以及 ${\mathcal S}'$）一般不是真正的线丛，而是线丛复幂的张量积；${\mathcal L}^w$ 只是记号，仅第一陈类 $w$ 有确切含义。其二，$w$ 住在 $H^2(\M(\alpha,p),\C)=(\Z\oplus\Lambda_\char)\otimes\C$，其中 $\Z$ 分量对应 $K_\M^{1/2}$、$\Lambda_\char\otimes\C$ 分量对应复化 Kähler 类；式~\eqref{eq:4.10} 的 $0\oplus(\cdots)$ 正是说第一分量取零。其三，$\eta$ 的负号不可漏：由 $[\,\omega_J/2\pi]=-\beta^*$ 与 $[B/2\pi]=\eta$ 及 $F+B=({\rm Im}\,\tau)\omega_J$ 得 $[F/2\pi]=-\eta$，脚注中的 ${\mathcal L}'={\mathcal L}^{-v}$ 与式~\eqref{eq:4.10} 中 $-\eta$ 是同一约定。其四，把式~\eqref{eq:4.2}、式~\eqref{eq:4.3} 代入可知 $\eta+i({\rm Im}\,\tau)\gamma^*=-({}^L\alpha-i{}^L\gamma)$，故 ${\mathcal L}^w$ 的指数恰为单值性的本征值 ${}^L\alpha-i{}^L\gamma$；这正是“线丛的指数等于单值性的本征值”这句总结的来源，也是核对本小节所有符号时最有效的检验。
\end{yizhu}


类比 \cite{ref8} 第 11.3 节对 $\Psi$ 或 $\theta$ 依赖的处理，我们将通过证明该结果在 $\gamma=0$、$\eta$ 等于格向量 $v\in \Lambda_\char$ 时成立来论证它。关键步骤是定义膜 ${\mathcal B}'$——粗略地说，它支撑在余切丛的零截面 $\M\subset T^*\M\subset \MH$ 上。一旦找到此膜，$({\mathcal B}_{c.c.},{\mathcal B}')$ 弦就给出由
$({\mathcal B}_{c.c.},{\mathcal B}_{c.c.})$ 弦所提供的环层的自然模层。通过辨明这个模，我们就能辨明该环层。

执行这一方案时会遇到一个重要细节。为在 $\eta\not=0$ 时定义膜
${\mathcal B}_{c.c}$（假定其 Chan-Paton 丛 ${\mathcal L}$ 平凡），需要如上所述取 $\beta^*=-({\rm Im}\,\tau)^{-1}\eta$。
这意味着从复结构 $I$ 的全纯观点看，$\MH$ 并非双有理于 $T^*\M$，而是双有理于其仿射形变，因此我们没有全纯嵌入 $\M\subset\MH$。
这初看会妨碍定义支撑在 $\M$ 上的膜 ${\mathcal B}'$。

然而，就目前目的而言，我们不关心复结构 $I$ 下全纯地发生什么。
我们关心的是 $K$ 型 $A$-模型，其中目标空间 $\MH$ 被理解为具有辛结构（为 $\omega_K$ 的某个倍数）的辛簇。
如第~\ref{sec:3.2}~节所述，作为具有辛结构 $\omega_K$ 的辛簇，$\MH$ 自然不依赖于 $\beta$。
这种 $\beta$-无关性是把 $\MH$ 看成复结构 $J$ 中平坦丛的模空间而得到的。

\begin{yizhu}[易错点]
「与 $\beta$ 无关」有精确含义：由式~\eqref{eq:3.77}，$\beta$ 只进入 $[\omega_J/2\pi]=0\oplus(-\beta^*)$，而 $K$ 型 $A$-模型用的是 $\omega_K$，其类 $[\omega_K/2\pi]=0\oplus(-\gamma^*)$ 不含 $\beta$。故移动 $\beta$ 不改变辛簇 $(\MH,\omega_K)$，只改变复结构 $I$ 下全纯嵌入是否存在（$\beta\ne0$ 时须取 $\beta^*=-({\rm Im}\,\tau)^{-1}\eta$ 才定义得出 c.c. 膜）。这正是第~\ref{sec:3.2}~节可以先把 $\MH$ 作为辛簇一次定下、再在它上面讨论与 $\beta$ 无关的 $A$-模型的原因；换言之，$\beta$ 破坏的是 $I$-全纯性，不是辛结构。
\end{yizhu}

当 $\beta\not=0$ 时，我们没有在复结构 $I$ 下全纯的嵌入 $\M\subset
\MH$，但确有相对于辛结构 $\omega_K$ 为拉格朗日的这种嵌入。
为得到这样的嵌入，从 $\beta=0$ 出发，用通常的、在复结构 $I$ 下全纯的嵌入 $\M\subset T^*\M\subset \MH$。按此定义，$\M$ 是复结构 $I$ 下的复拉格朗日子流形——即在复结构 $I$ 下全纯，且对全纯辛形式 $\Omega_I=\omega_J+i\omega_K$ 为拉格朗日。特别地，$\M$ 对 $\omega_K$ 为拉格朗日。

现在转换观点，把 $\MH$ 看成复结构为 $J$ 的复辛流形。由此观点，$\MH$ 典范地不依赖于 $\beta$。
（改变 $\beta$ 时不会遇到奇点，因为我们取 $\alpha$ 正则，以求得涉及 $\M(\alpha,p)$ 上 ${\mathcal D}$-模的答案。）
改变 $\beta$ 时，我们不改变全纯辛结构
$\Omega_J=\omega_K+i\omega_I$。特别地，我们在保持辛结构 $\omega_K$ 固定的情况下改变 $\beta$。因此上一段找到的子流形 $\M\subset \MH$，若仅看成 $K$ 型拉格朗日子流形，则对任意
$\beta$ 都有自然定义。若改变 $\beta$ 时保持由 $\MH$ 中一点所决定的（复结构 $J$ 下的）带分歧平坦丛不变，我们就取它逐字地不依赖于 $\beta$。

为得到支撑在 $\M$ 上的 $A$-膜 ${\mathcal B}'$，需要一个酉线丛 ${\mathcal L}'\to \M$
使其曲率 $F$ 满足 $F+B|_{\mathcal M}=0$。${\mathcal L}'$ 将成为膜 ${\mathcal B}'$ 的 Chan-Paton
线丛。这里 $B$ 是 $B$-场，由于我们假定 $\theta=0$，它由 $\eta$ 决定。事实上，$B$ 是 $\MH$ 上的闭二形式，其上同调类为 $0\oplus \eta\subset
(\Z\oplus\Lambda_\char)\otimes_\Z\R=H^2(\M(\alpha,
p);\R)$。对一般的 $\eta$，该上同调
类不是格向量，因此不存在合适的线丛
${\mathcal L}'$ 或膜 ${\mathcal B}'$。
然而，若对某个 $v\in \Lambda_\char$ 有 $\eta=v$，则可取
\footnote{下式中指数的负号来源如下。$B$ 与 $\eta$ 的关系中没有负号：$[B/2\pi]=\eta$。
故由 $[F+B]=0$ 得 $[F/2\pi]=-\eta$。因此若 $\eta=v$，Chan-Paton 线丛的第一陈类为 $-v$。}
${\mathcal L}'={\mathcal L}^{-v}$，
即 ${\mathcal L}'$ 是第一陈类为
$0\oplus (-v)$ 的线丛。

如此定义了 $K$ 型 $A$-膜 ${\mathcal B}'$ 后，考虑 $({\mathcal B}_{c.c.},
{\mathcal B}')$ 弦；它们将给出由 $({\mathcal B}_{c.c.},{\mathcal B}_{c.c.})$ 弦导出的环层的模层。在量子化
$({\mathcal B}_{c.c.},{\mathcal B}')$ 弦时，与 $\eta=0$ 情形相比唯一的新之处是：膜 ${\mathcal B}'$ 的 Chan-Paton
线丛不再是平凡的，而是
${\mathcal L}'={\mathcal L}^{-v}$。这意味着 $({\mathcal B}_{c.c.},{\mathcal B}')$ 弦的层不再是
$K_\M^{1/2}$ 的截面层，而是
$K_\M^{1/2}\otimes {\mathcal L}^{-v}$ 的截面层。这表明式~\eqref{eq:4.10} 对 $\gamma=0$、$\eta
\in\Lambda_\char$ 成立，因而对一切 $\gamma,\eta$ 成立。

现在可以用 ${\mathcal D}$-模而非镜像对称，重新表述带驯服分歧的平坦丛的几何朗兰兹对偶。
在几何朗兰兹对应的左侧，考虑
结构群为 $^L\neg G$ 的带分歧平坦丛 $E\to C$，其单值性所对应的轨道闭包中含有半单元素
$U=\exp(-2\pi({}^L\neg\alpha-i{}^L\neg\gamma))$。右侧，这种带分歧平坦丛的对偶是
$\M(\alpha,p;G)$ 上被 $K_\M^{1/2}\otimes
{\mathcal L}^{-(\eta+i({\rm Im}\,\tau)\gamma^*)}$ 扭曲的微分算子层的模层，而它等同于
$K_\M^{1/2}\otimes
{\mathcal L}^{^L\neg\alpha-i{}^L\neg\gamma}$。于是，“线丛的指数”与“单值性的本征值”相同。
$S$-对偶在两类对象之间建立自然的对应。
（对具有非正则单值性的平坦丛，此陈述需要一些澄清，我们留到第~\ref{sec:4.6}~节。）
这一版本的几何朗兰兹对偶已在数学上被猜想。参见 \cite{ref7} 第 9.4 节的综述以及引言中的其他文献。

\bigskip\noindent{\it 非正式的说明}

我们刚才给出的说明是我们所知的最精确的一个。
然而，有些读者可能更喜欢另一种解释，我们在此仅作非正式的介绍。

考虑任一复簇 ${\mathcal M}$，令 $X=T^*{\mathcal M}$ 带上其自然全纯辛形式 $\Omega_X$。$X$ 的量子化给出 ${\mathcal M}$ 上被 $K_\M^{1/2}$ 扭曲的微分算子。现在把 $X$ 换成一个作为余切丛仿射形变的簇 $Y$，即存在全纯纤维化 $\pi:Y\to {\mathcal M}$，其纤维与余切丛的纤维相同，但 $\pi$ 不容许全纯截面。局部上可以取一个全纯截面并把 $Y$ 与 $X$ 等同；整体上，此等同的障碍由元素 $w\in H^1({\mathcal M},T^*{\mathcal M})$ 决定。进而要求 $Y$ 容许全纯辛形式 $\Omega_Y$，它在局部（取 $\pi$ 的局部全纯截面、从而把 $Y$ 局部等同于 $X$ 之后）与 $\Omega_X$ 重合。这等价于说 $w$ 可由 ${\mathcal M}$ 上一个复值 $(1,1)$-形式表示，该形式被 $\bar\partial$ 与 $\partial$ 算子零化，从而特别地被 $d=\partial+\bar\partial$ 零化。若 ${\mathcal M}$ 是紧 Kähler 流形，此条件自动满足。
在此情形下，自然可以断言：以辛形式 $\Omega_Y$ 量子化 $Y$，给出 ${\mathcal M}$ 上被 $K_{\mathcal M}^{1/2}\otimes {\mathcal L}^w$ 扭曲的微分算子，其中 ${\mathcal L}^w$ 象征性地表示第一陈类为 $w$ 的“线丛”。

刻画扭曲参数 $w$ 的一个便利方式如下。全纯
辛形式 $\Omega_X$ 的上同调类为零。事实上，若 $q^\alpha$ 为 ${\mathcal M}$ 上的局部坐标、$p_\alpha$ 为相应的正则动量，则 $\Omega_X=\sum_\alpha
\,dp_\alpha\wedge dq^\alpha=d(\sum_\alpha\,p_\alpha\,dq^\alpha)$。
这里的一形式 $\lambda=
\sum_\alpha \,p_\alpha\,dq^\alpha$ 整体有定义，故 $[\Omega_X]=0$。相反，一个标准结果是：$\Omega_Y$ 的上同调类等于闭 $(1,1)$-形式 $w$ 在
$H^2({\mathcal M},\mathbb{C})$ 中的类（更确切地说，是该类拉回到 $Y$
上的类）。我们就直接把这个上同调类记为 $w$。

\begin{yizhu}[含糊点拨]
同一个 $w$ 在此有两副面孔：本节开头它是 $H^1(\M,T^*\M)$ 中的障碍类（$\pi:Y\to\M$ 无全纯截面时的过渡余循环），这里又是 $H^2(\M,\mathbb{C})$ 中闭 $(1,1)$-形式的类。两者一致，是因为 $T^*\M$ 的全纯截面层就是 $\Omega^1_\M$，故 $H^1(\M,T^*\M)=H^1(\M,\Omega^1)$；而 $w$ 被 $\partial$ 与 $\bar\partial$ 同时零化，给出 $(1,1)$-型且 $\bar\partial$-闭的代表，于是 Hodge 分解把 $H^1(\M,\Omega^1)$ 与 $H^{1,1}(\M)$ 等同。所以「$\Omega_Y$ 的上同调类等于 $w$ 的类（拉回到 $Y$）」与障碍类的说法是同一件事的两种表述。
\end{yizhu}

现在回到我们的问题：$\MH$ 在复结构 $I$ 下（除高余维集外）是 $\M$ 的余切丛的仿射形变，
全纯辛形式为
$\Omega=({\rm Im}\,\tau)\Omega_I=({\rm Im}\,\tau)(\omega_J+i\omega_K)$。
由式~\eqref{eq:3.77}，$\Omega$ 的上同调类为
$w=-({\rm Im}\,\tau)(\beta^*+i\gamma^*)$。为构造把 $A$-膜范畴与 ${\mathcal D}$-模联系起来所需的 c.c. 膜，我们不得不取 $\beta^*=-({\rm Im}\,\tau)^{-1}\eta$，故扭曲参数实际上是 $w=\eta-i({\rm Im}\,\tau)\,\gamma^*$，与上面所得一致。

\begin{yizhu}[原文疑误]
此处（PDF 第 96 页）的启发式推导与式~\eqref{eq:4.10} 相差一个符号，疑有笔误。把 $\beta^*=-({\rm Im}\,\tau)^{-1}\eta$ 代入本页所给 $w=-({\rm Im}\,\tau)(\beta^*+i\gamma^*)$，得 $w=+\eta-i({\rm Im}\,\tau)\gamma^*$，即本页印出的结果；而式~\eqref{eq:4.10} 为 $w=0\oplus(-\eta-i({\rm Im}\,\tau)\gamma^*)$。二者必有其一有误。倾向于认为本页（或式 (4.10) 的 $\eta$ 符号）有笔误，理由是：（1）上文关于 ${\mathcal L}'={\mathcal L}^{-v}$ 的脚注明确由 $[F+B]=0$ 与 $[B/2\pi]=\eta$ 得 $[F/2\pi]=-\eta$，并据此取 ${\mathcal L}'={\mathcal L}^{-v}$，这与式 (4.10) 的 $-\eta$ 自洽；（2）本小节末“被 $K_\M^{1/2}\otimes{\mathcal L}^{-(\eta+i({\rm Im}\tau)\gamma^*)}$ 扭曲”以及“线丛的指数等于单值性的本征值”这两处表述都要求 $w$ 中 $\eta$ 带负号。照原文保留公式未作改动，建议核对 PDF 第 96 页该句与式 (4.10) 的符号。
\end{yizhu}

我们认为这一解释只是启发式的（尽管对某些读者也许更易理解），原因在于我们无权任意假定 c.c. 膜应具有何种性质；这应当作为量子场论标准框架的一部分被推导出来。上述论证连同 \cite{ref8} 第 11 节的论证，正是为此所作的尝试。

\bigskip\noindent{\it 格平移对称性}


有人现在会问：在上述推导中为何没有看到按格向量平移 $\eta$ 的对称性，例如对 $v_0\in\Lambda_\char$ 的 $\eta\to\eta+v_0$？
事实上，我们是在假定 c.c. 膜的 Chan-Paton 丛 ${\mathcal L}$ 平凡时丢掉这一对称性的。格平移 $\eta\to\eta+v_0$ 作用在所有膜上的方式是：把它们的 Chan-Paton 丛张量上 ${\mathcal L}^{-v_0}$。要恢复该对称性，只需补充额外的余迷向膜 ${\mathcal B}_{c.c.}^{v_0}$，它们由
${\mathcal B}_{c.c.}$ 把 $\eta$ 平移 $v_0$ 并把 Chan-Paton 丛张量上 ${\mathcal L}^{-v_0}$ 而得到。
我们不需要再引入任何 ${\mathcal B}'$ 膜。
事实上，上面构造出的全部 ${\mathcal B}'$ 膜已具有该平移对称性，因为这些膜是对一切可能的格向量 $v$ 与一切 $\beta$ 定义的。

从数学观点看，考察这一平移对称性所得到的东西叫做 Morita 等价。
若 ${\mathcal L}^{v_0}$ 是真正的线丛（而非线丛的复幂），则作用在
$K_\M^{1/2}\otimes {\mathcal L}^w\otimes {\mathcal L}^{v_0}$ 截面上的微分算子层 Morita 等价于作用在
$K_\M^{1/2}\otimes {\mathcal L}^w$ 截面上的微分算子层。该 Morita 等价
由考察一个“双模”来建立，在目前语境下这个双模就是
$({\mathcal B}_{c.c.},{\mathcal B}_{c.c.}^{v_0})$ 弦的层。

\begin{yizhu}[含糊点拨]
Morita 等价把 $\eta\to\eta+v_0$ 的平移实现为环层的模范畴变换：$\cD_{K_\M^{1/2}\otimes\cL^w}$ 与 $\cD_{K_\M^{1/2}\otimes\cL^w\otimes\cL^{v_0}}$ 的模范畴等价，实现它的双模正是 $(\CB_\cc,\CB_\cc^{v_0})$ 弦层。条件 $v_0\in\Lambda_\char$ 不可省：只有格向量才给出真正的线丛 $\cL^{v_0}$，否则只是线丛的复幂，连 $\cD$ 作为环层都无从定义——这正是正文要加「若 $\cL^{v_0}$ 是真正的线丛」这一限制的原因。于是物理上的格平移对称性，在 $\cD$-模一侧恰好是「$w$ 只在模 $\Lambda_\char$ 的意义下有内容」这一事实。
\end{yizhu}

\bigskip\noindent{\it $\mathbb{L}$ 型情形的类比}

我们所说的一切，对任意列维子群 $\mathbb{L}$（而非此前考虑的 $\mathbb{L}=\TT$）都有直接类比。
记 $\mathbb{L}$ 的外尔群（Weyl group）为 ${\mathcal W}_\mathbb{L}$。

对一般的 $\mathbb{L}$，参数 $(\alpha,\beta,\gamma,\eta)$ 被限制为
${\mathcal W}_\mathbb{L}$-不变量。在此理解下，把对偶表述为 $\Psi=\infty$ 处 $B$-模型与 $\Psi=0$ 处 $A$-模型之间的关系，无须实质修改。该对偶在
$\M_{H,{}^L\neg\mathbb{L}}({}^L\neg\alpha,{}^L\neg\beta,{}^L\neg\gamma,{}^L\neg\eta;p,{}^L\neg G)$
的 $B$-膜与 $\M_{H,\mathbb{L}}(\alpha,\beta,\gamma,\eta;p,G)$ 的 $A$-膜之间建立自然对应，参数之间照例满足同样的关系。

要用 ${\mathcal D}$-模表述对偶，必须假定 $\alpha$ 是
$\mathbb{L}$-正则的。此时记由对 $(\mathbb{L},\alpha)$ 决定的抛物子群为 $\EUP$，并引入 $\M_\mathbb{L}(\alpha;p,G)$，即 $C$ 上在点 $p$ 处具有 $\EUP$ 型抛物结构的 $G$-丛的模空间。当 $\beta=\gamma=0$ 时，如式~\eqref{eq:3.73} 所见，希格斯丛模空间\footnote{此处省略 $\eta$，因为它不进入经典
几何。} $\M_{H,\mathbb{L}}(\alpha,\beta,\gamma;p,G)$ 以
$T^*{\M}_\mathbb{L}(\alpha;p,G)$
为稠密开子空间。当 $\beta,\gamma\not=0$ 时，
$\M_{H,\mathbb{L}}(\alpha,\beta,\gamma;p,G)$ 以余切丛的一个仿射形变为稠密开子空间。

\begin{yizhu}[易错点]
$\mathbb{L}$-正则指 $\alpha$ 在外尔群 ${\mathcal W}$ 中的稳定化子恰为 ${\mathcal W}_\mathbb{L}$。只有此时由 $(\mathbb{L},\alpha)$ 才能唯一确定以 $\mathbb{L}$ 为列维的抛物子群 $\EUP$，$\M_\mathbb{L}(\alpha;p,G)$ 也才良定义，式~\eqref{eq:3.73} 的稠密开嵌入 $T^*\M_\mathbb{L}\subset\M_{H,\mathbb{L}}$ 才成立。若 $\alpha$ 不是 $\mathbb{L}$-正则的，就必须换用更大的列维群 $\mathbb{L}'\supsetneq\mathbb{L}$——这正是第~\ref{sec:3.6}~节讨论分歧模空间奇点时的做法（$\alpha$ 趋于非 $\mathbb{L}$-正则值时，纤维是 $\mathbb{L}'/\mathbb{L}$ 型的消没闭链）。
\end{yizhu}

这些几何事实连同典范余迷向膜的存在性，可以完全照上述方式用来得到自然映射：从 $A$-膜到 ${\mathcal D}$-模，
其中 $A$-膜取在 $\M_{H,\mathbb{L}}(\alpha,\beta,\gamma,\eta;p,G)$ 上。
所有关键陈述与论证都相同。唯一的改变是：陈述现在涉及 $\M_{H,\mathbb{L}}$ 与 $\M_\mathbb{L}$（而非
$\MH$ 与 $\M$），且参数是 ${\mathcal W}_\mathbb{L}$-不变量。

在某种意义上，对一般 $\mathbb{L}$ 的陈述比
$\mathbb{L}=\TT$ 的陈述略欠自然。
$\mathbb{L}$-正则的 $\alpha$ 的不同选取（当然对应于
$^L\neg\eta$ 的不同选取）导致
$\EUP$ 的不同选取。因此它们导致不同模空间上的 ${\mathcal D}$-模，
而这些模空间都对应着具有不同类型抛物结构的丛。
抛物结构为 $\EUP$ 型，
其中 $\EUP$ 可取为 $G_\C$ 的任一以 $\mathbb{L}$ 为列维子群的抛物子群。
不过，这些不同的空间彼此双有理等价。
